\chapter[Steppe, Forest, and Ocean Histories]{People of the Steppe, Forest, and Ocean Make History Too}
\section{A Rope Connecting an Ocean}
Pick up something neither bronze, coin, nor book: a brown-black rope, slightly thicker than your thumb, stiff and rough, twisted fibre by fibre from coconut husks.

Pacific islanders use it for boatbuilding. Polynesian and Micronesian voyaging canoes were not nailed together; many islands lacked iron ore. Carpenters fitted planks, pierced their edges, sewed and lashed them with coir, and waterproofed seams with resin and crushed plant fibres. A fifteen- or twenty-metre double canoe with outrigger supports depended on thousands of lashings from keel to mast. Broken rope meant a broken ship.

Such vessels carried people, pigs, chickens, dogs, taro shoots, banana stems, and breadfruit cuttings across Earth's largest waters. From east of today's Philippines, voyagers over millennia reached Tonga, Samoa, Tahiti, Hawai'i, Easter Island, and finally New Zealand, while European sailors still fearfully hugged coasts without losing sight of land.

Look again at this unlettered, undecorated rope destined for a museum corner drawer. Without it, planks remain timber and islands isolated. A third of Earth's surface was tied together with rope.

Why do history books almost always begin with walls, palaces, and farmers? Boatbuilders, riders, reindeer followers, and camel drivers also made history, yet appear mainly as invaders or background. Follow ropes, horses, and caravans beyond the book written by settlers.

\section{A Day at the Horse Market}
Enter a scene.

Time: autumn in the late 1570s, early in Ming Emperor Wanli's reign. Place: an officially designated horse market on a riverbank beyond a Great Wall pass. Several years after the Longqing peace agreement, Ming and the Mongol right-wing groups have substituted markets for battlefields. This is one agreed market.

At dawn, frost remains on grass. Mongol herder Batur arrives with a dozen horses and sheep after three days from northern pasture. He wears an old sheepskin coat and a belt knife with a polished-worn sheath. Horse snorts, sheep smells, and steaming human breath mingle in cold air.

Across the riverbank, Han traders from Xuanhua, Datong, and other towns erect stalls: stone-hard tea bricks, blue cloth, iron pots, needles, scissors, silk thread, and grain. Interpreters dressed in two styles switch rapidly between Mongolian and Chinese, their sharp voices enabling bargains.

Batur examines a pot, turning it and tapping its bottom for sound. The steppe lacks iron and smiths; a good pot is cookware, dowry, bride-gift, and durable inheritable wealth. He touches tea bricks too. Meat-and-milk diets welcome tea to cut richness and aid digestion, and habitual drinkers miss it daily. His wife also requested cloth and needles.

The Han bazong, the officer keeping market order, watches Batur's horses. Ming cavalry and post stations need them; inland farms need animals too. Wind-and-snow pasture raises endurance beyond that of penned horses. Inland workshops also await hides and wool.

By afternoon the bargain is done: horses for tea, cloth, and a pot; sheep for grain and ironware. Batur lashes the pot behind his saddle and turns north towards tent, seasonal routes, water, and pasture. Traders pack and head south towards walls, villages, fields, and ancestral halls.

They return in opposite directions, neither wishing to become the other. Yet neither household can do without the other.

\section{People Textbooks Cannot Accommodate}
Before praising the exchange, put an uncomfortable question on the table.

Recall textbook structure: farming brings surplus and villages, then cities, writing, law, and states. Civilisation appears and history officially begins. Agriculture, houses, and settlement mark progress; steppe herders, forest hunters, and voyagers lack a station on the line. They enter as sudden barbarian raiders who vanish after plunder, or uncivilised peoples waiting at page margins for development and incorporation.

Three gaps undermine this account.

First, it cannot explain the market. If steppe people only plunder, why travel three days to trade? If Batur lacks culture, who taught him to test pots, and why do officials eagerly await his horses?

Second, it cannot explain oceans. Centuries before European oceanic voyages, Pacific islanders travelled thousands of kilometres in nailless boats. If settlement is civilisation's prerequisite, what about people making ships home and seas roads? Are stars, swells, and monsoons not knowledge?

Third, it quietly treats mobility as transition: pastoralists have not learned farming, hunters housebuilding, sailors where to settle. Mobile lives become unfinished settled drafts awaiting their end. Test this directly: \textbf{is mobility failed settlement or another complete, mature way of life?}

Choose provisionally: if generations remain on the move, is that deficiency or choice?

\section{What the Steppe Thinks and the Sea Says}
Most surviving ancient writing about steppe people comes from farmers, unsurprisingly given urban preservation. But strengthen settlers' position first: it contains real reasoning, not only prejudice.

Agrarian states run on taxes, requiring countable land, people, and cattle registered for autumn collection. A mobile herder is difficult to tax, recruit, or govern. Administrative difficulty precedes cultural arrogance. Next comes real fear: rapid cavalry did raid villages; archival bloodstains are not invented. Then comes prejudice: urban writers measure others by their own lives, describing absent walls and ritual. Holding the pen, they still dominate what we read three millennia later. This does not make every source a lie; \textbf{surviving voices are not all voices}.

What does the steppe say? Less survives, but some statements carry enormous weight.

Early-eighth-century Turkic monuments in Mongolia's Orkhon Valley record the Second Turkic Khaganate in Turkic and Chinese. The best-known honours Kül Tegin. Its Turkic text gives a khagan's warning, broadly: the powerful southern Tang state speaks sweetly and offers soft gifts, attracting distant peoples. Once they draw near and settle, harmful designs grow and their people gradually disappear. Remain on your steppe and preserve independence rather than settling there.

Consider this thirteen-hundred-year-old steppe declaration. It offers mature political analysis, not textbook barbarian ignorance: power creates dependence through inducements; a community faces wealthy neighbours while defending itself. Its author understood settled civilisation's force and therefore declined entry. Settlement was not beyond his ability but against his choice.

Move to the ocean for a third voice. In 1769, James Cook's expedition reached Tahiti and took aboard local priest-navigator Tupaia. Asked about nearby islands, he recalled dozens of names, directions, distances, and sailing days across thousands of kilometres. Expedition scholars helped make a chart. Europe long dismissed its unfamiliar logic as Indigenous doodling, until comparison revealed strikingly accurate island relationships in another coordinate system: an ocean civilisation's world map.

The story raises a troubling question. Europeans recorded Cook discovering islands while someone beside him could recite their routes. Who discovered whom? In history, discovery often means respectably saying that our people first heard of something.

Settlers, steppe, and sea each speak from interests and positions, not pure facts alone. Hearing that position matters more than memorising every statement.

\section[Thinking Bridge: Mutual Dependence]{Thinking Bridge: List Who Cannot Do without Whom}
Whenever civilisation and barbarians appear, use a portable list of mutual dependence, four columns and a connecting line.

\begin{itemize}
\item \textbf{What we have}: local abundance available for exchange.
\item \textbf{What we lack}: things difficult to make or grow but hard to live without.
\item \textbf{What they have and lack}: complete the same columns from their side.
\item \textbf{The connecting route}: trading road, horse market, or sea lane joining surpluses and needs.
\end{itemize}
At the market, the steppe offers horses, sheep, hides, and wool, needing iron, tea, cloth, and grain. Farming regions reverse that list, connected by frontier markets and caravan roads. The civilisation-barbarian frame collapses into two incomplete, interlocking worlds. Ming needed Mongol horses no less than Mongols needed Ming pots.

At sea, big islands offer stone, timber, and farmland; coral atolls offer coconuts and fish and may import even obsidian for sharpening from hundreds of kilometres away. Foreign-source stone tools on small islands establish lasting exchange. Navigation was not adventurers' entertainment but residents' lifeline. Ships were roads; builders and navigators became hubs.

Today, cities offer wages, schools, and hospitals but need food and labour; villages offer land and produce but lack income and services. Roads, railways, and migrant workers connect them. Or compare you with your delivery rider. Four columns reveal dependence immediately.

But the list has limits. It corrects imagined one-way relations without establishing fairness. Mutual need does \textbf{not} mean equality: coercive bargains, dominant sea powers, and algorithm-set delivery rates persist. Who depends on whom differs from who exploits whom, requiring another chapter's tools. Dismantling assumed superiority already repays the effort.

\section{Sima Qian's Children Riding Sheep}
Read an early systematic Chinese description of steppe life, Sima Qian's Records of the Grand Historian, Account of the Xiongnu, completed over twenty-one centuries ago. Without visiting the steppe, he gathered a century's diplomatic archives, reports, and frontier accounts:

\begin{quote}
They migrate after water and grass, without walled cities, fixed residences, or farming occupations, yet each has allotted land.
\end{quote}
\begin{quote}
Children able to ride sheep draw bows at birds and mice; growing older, they shoot foxes and hares for food.
\end{quote}
The first describes mobile life without fixed agricultural settlement but with divided pasture. The second shows children beginning on sheep, hunting small creatures, then larger prey for food.

Read three evidential layers.

First, what did he recognise? Precise adaptation. Low, slower sheep teach children riding; small targets train sight and aim; food combines play, education, and production. Preparation for adult mounted archery begins with walking. This is a school without classrooms, not idle wandering. He also recorded offices beneath the chanyu, including Left and Right Worthy Kings and Guli Kings, each with territories. He knew an institutional political body, not a mob.

Second, where is his frame? His phrase begins with \textbf{without}: walls, fixed homes, agriculture. Counting absences establishes settled reference points, like introducing your deskmate by lacking merit certificates or maths classes. Even later praise cannot undo that coordinate system. \textbf{A text describes both a world and its writer's position}. Read both or evidence can mislead.

Third, what can it prove? Han observers recognised steppe organisation, skills, and logic; the four Chinese characters for allotted land show defined pasture rather than random wandering. It cannot prove innate barbarity or belligerence: neither quotation mentions plunder. That impression comes from warfare accounts agrarian dynasties particularly favoured. Unequal documentary quantities leave fingerprints of perspective.

\section{Why Did Cavalry Ride South?}
Evidence and tools allow our largest factual problem: for over two millennia, Xiongnu, Turks, Uyghurs, and Mongols repeatedly pressed settled dynasties from the steppe. \textbf{Why?} Three causal explanations do not justify war. Explaining mechanisms is not praising fighting.

\textbf{Explanation one: disaster and survival.}

Steppe livelihoods depend on harsher weather than farming. A white disaster, deep snow blocking forage, can kill hundreds of thousands of animals in one winter. Livestock are food, savings, and transport; death bankrupts households. Raiding south then means survival. Vulnerability and some disaster-raid timing support this. Yet many major campaigns occurred during prosperity, with well-fed herds; disasters weakened horses too much for large wars. Survival also cannot explain remaining to establish decades-long regimes rather than taking grain home.

\textbf{Explanation two: blocked commerce turns into war.}

Sixteenth-century Altan Khan repeatedly petitioned Ming for markets, returning captives and punishing raiders. Refusals continued, sometimes with envoys killed. In 1550 he reached Beijing in the Gengxu Incident, still demanding horse markets. The 1571 Longqing agreement finally opened tribute and trade, bringing decades of broad frontier stability and our opening market. Raiders often sought exactly what closed markets denied: iron, grain, tea, cloth. Why risk life if goods can be bought? Yet this cannot explain every age: early Han already practised marriage diplomacy and frontier trade. Reducing war to failed business also erases internal steppe politics, making khagans merchants denied licences.

\textbf{Explanation three: wealth holds political networks together.}

Steppe coalitions need distributions of silk, precious metals, tea, and iron to chiefs. Followers expect rewards; unrewarded alliances dissolve. Pastures produce horses and hides, not these valuables. Three routes provide them: trade, gifts through tributary forms of exchange, or plunder. Successive rulers calculated among these options. Genghis Khan and successors extended the system by protecting trade itself: Eurasian relay stations and official tablets supported merchants travelling from the Black Sea to Dadu. Some historians describe steppe empires as shadows of agrarian worlds, half their political building materials borrowed from neighbours. This unites raiding and trade-protecting khagans, but erases ordinary herders: why soldiers and horse-tending children followed, and what they gained, remains unclear.

Disaster explains some timing, trade some targets, political networks some scale. \textbf{One why changes with the camera's position: pasture, market, or royal camp}. All three show reality; no single frame contains the whole.

\section[Further Challenge: Dismantling Progress]{Further Challenge: Dismantle the Ladder of Progress}
Now confront the theoretical ladder hidden in countless textbooks.

Lewis Henry Morgan's Ancient Society, published in 1877, ranked humanity through savagery, barbarism, and civilisation: gathering and hunting below, pastoralism and farming between, cities and writing above. For over a century this framework influenced textbooks worldwide, implying mobile lives had not yet climbed to settlement, everyone's final station.

Evidence, not emotion, has largely dismantled it in three forms.

\textbf{First: chronology.} If pastoralism were agriculture's childhood, it should precede it. Instead villages and domestic animals came first, mature mounted nomadism much later with developed riding. Herd management, rotational pasture, seasons, breeding, and food preservation form a specialised system, not failed farming. Pastoralism is agriculture's sibling: one stays, one turns otherwise unusable grasslands into milk, meat, wool, hides, and mobility. That is division of labour, not repeating a school year.

\textbf{Second: real frontiers.} Twentieth-century scholar Owen Lattimore saw walls as interfaces, not lines dividing civilisation and barbarism. Agrarian states developed walls, markets, marriage diplomacy, and rewards through steppe relations; steppe coalitions developed through negotiating with them. Exchange, intermarriage, flight, and allegiance crossed boundaries. A one-sided history of civilisation must first pretend that interface absent.

\textbf{Third: mobile knowledge itself.} This strikes hardest. Examine several settings.

Pacific navigators lacked magnetic compasses, sextants, and paper charts but possessed mental star compasses, fixed rising and setting directions, and eyes reading swells bent by islands tens of kilometres beyond sight. Birds and lagoon-green reflections under clouds added clues. Marshall Islanders made stick-and-shell charts representing swells as shore-based teaching aids. These supported extraordinary expansion long dismissed in Europe as accidental drift. In 1976, Hawaiians built the traditional double canoe Hōkūle'a and invited Satawal navigator Mau Piailug to sail without modern instruments roughly four thousand kilometres to Tahiti. Drift theory failed before demonstrated knowledge written in stars and waves rather than paper.

In the Greater Khingan forests, reindeer-herding Evenki have moved for centuries, reading snow crust, moss, and hoofprints to choose tomorrow's route. Outsiders see uniform forest; residents see signs. Amazonian fruit-tree concentrations once called untouched wilderness increasingly appear deliberately planted and managed over generations. Forests were cultivated homes, not merely places passed through.

In the Sahara, Tuareg caravans have moved salt for over a millennium from mines to waiting oases and cities. Stars, winds, dunes, even sand's smell and colour guide hundreds of kilometres without signposts. Salt returns as grain and cloth: barren desert becomes a road surface.

Finally, a polar warning against misjudgement. In 1845, Franklin's two advanced British naval vessels entered the Arctic seeking the Northwest Passage and vanished, leaving no survivors among over a hundred crew. Later searchers repeatedly heard clear Inuit accounts of ships and survivors abandoning them southwards. Many nineteenth-century Britons dismissed the testimony or blamed Inuit for the catastrophe. Wreck discoveries in 2014 and 2016 lay near areas preserved in oral memory. Advanced instruments had lost to the memories of people deemed unknowing. Remember these ships before dismissing supposedly backward knowledge.

Together the evidence makes progress's ladder a nineteenth-century invention, not destiny. \textbf{Mobility is not settlement's unfinished tense}. Horsebacks, camels, canoes, and forest snow trails sustain complete worlds, schools, engineers, and archives, inscribed on stones, sung in songs, or tied into coconut-rope knots.

\section{Today: People Still Moving}
These questions live in news and neighbourhoods, not museums.

After Hōkūle'a, traditional star navigation revived across Tahiti, Samoa, New Zealand, the Cook Islands, and elsewhere. Communities rebuilt canoes, relearned ancestral star maps, and voyaged around the Pacific; some schools added navigation lessons. This is more than nostalgia: supposedly dead knowledge still crosses oceans. Tradition is a technology ready to sail again, not a glass-case specimen.

Pastoral regions face conflicting pressures. Settlements bring genuinely desired schools, hospitals, heat, and electricity. Concentration can also overload pasture, prevent seasonal rotation, and break transmission of elders' snow, grass, and water knowledge. In Mongolia, disasters and livelihood troubles drive herders into Ulaanbaatar's expanding ger districts, still in tents but without pastures or secure jobs. Settled administratively, perhaps, but not in life. Who judges improvement: tidy planning diagrams or someone knowing every spring-flooded hollow?

Cities depend on moving people: delivery riders, lorry drivers, construction workers, migrant parents, digital nomads with suitcases. Every punctual meal arrives through someone on the road. Yet society praises homeownership, registration, and settling down. Annual family questions reactivate the nineteenth-century ladder: mobile workers support everything while their identity remains temporary, awaiting confirmation. Ancient pastoralists described without fixed dwellings and modern riders registered without permanent addresses inhabit the same grammar.

Return once more to coir. As islanders revive voyaging, rising seas threaten ancestral atolls. Those most skilled in surviving through movement may one day teach a final lesson to a world assuming it could stay still forever.

\section[Your Question: Must We Settle?]{Your Question: Is Settlement the Only Answer?}
Return to the horse market at dusk.

Batur lashes the pot to his saddle and heads home, declining urban life not because he cannot farm but because his world is complete: sheepback schools, livestock banks, seasonal roads, epic and stone archives. The trader heading south likewise declines his life. Their worlds meet, exchange necessities, and return home.

We have spoken for them at length. The final question is yours:

\textbf{If a life requires continued movement---herding, sailing, construction work, travel among islands---why assume settlement brings happiness rather than deprivation?}

Weigh both sides first.

On one side, schools, hospitals, and heat are real benefits; stability can be as sincerely desired as mobility. Romanticising pastoral or seafaring freedom from a heated room is another arrogance. Winter kills herds and boats capsize. Nobody owes you a life serving your fantasy of freedom.

On the other, settlement imposed for people's own good has destroyed complete worlds. Refusal of your designed good life may reflect knowledge you lack, like the khagan's awareness of sweet words and soft gifts or Inuit elders' remembered seas dismissed by the Royal Navy.

Is settlement the only answer? Who should judge a life: planners or its inhabitants? Answer with reasons. Your home, school, and familiar stability have already cast a quiet vote. Seeing that vote is where your answer begins.
